\documentclass[11pt,a4paper]{article}

% Packages
\usepackage[utf8]{inputenc}
\usepackage[T1]{fontenc}
\usepackage[margin=0.7in]{geometry}
\usepackage{enumitem}
\usepackage{hyperref}
\usepackage{titlesec}
\usepackage{xcolor}

% Formatting
\pagestyle{empty}
\setlength{\parindent}{0pt}
\setlist[itemize]{leftmargin=12pt, noitemsep, topsep=2pt}

% Section formatting
\titleformat{\section}{\Large\bfseries}{}{0em}{}[\titlerule]
\titlespacing*{\section}{0pt}{8pt}{6pt}

\titleformat{\subsection}[runin]{\bfseries}{}{0em}{}
\titlespacing*{\subsection}{0pt}{6pt}{0pt}

% Colors
\definecolor{primaryblue}{RGB}{0,82,155}
\definecolor{accentblue}{RGB}{0,120,215}
\definecolor{mediumgray}{RGB}{100,100,100}

% Hyperlinks
\hypersetup{
    colorlinks=true,
    urlcolor=accentblue
}

% Custom commands
\newcommand{\highlight}[1]{\textcolor{primaryblue}{\textbf{#1}}}
\newcommand{\keypoint}[1]{\textbf{#1}}
\newcommand{\skilltitle}[1]{\normalsize\textcolor{primaryblue}{\textbf{#1}}}

\begin{document}

\begin{center}
    {\LARGE \textbf{ERIC MAYEUX}}\\
    \vspace{2pt}
    {\normalsize Scrum Master — Agile / SAFe}\\
    \vspace{3pt}
    43 rue Levis, Longueuil J4H 4B6, QC, Canada\\
    \href{mailto:mayeux.e@gmail.com}{mayeux.e@gmail.com} | 438-884-7645
\end{center}

\vspace{6pt}

\section*{Profil Professionnel}
\noindent
\highlight{Scrum Master Certifié} avec \keypoint{10+ ans d'expérience} en \highlight{méthodologie agile} et \highlight{facilitation} de cérémonies pour plusieurs squads. Reconnu pour ma \highlight{communication} efficace, mon \highlight{écoute active} et ma capacité à favoriser la \highlight{collaboration} et l'efficacité d'équipe. Expert en \highlight{Framework Scrum} et \highlight{SAFe} avec une approche \highlight{orientée résultats}. Ambassadeur des comportements BN : \highlight{Agilité}, Complicité et Pouvoir d'Agir. Démontre une forte \highlight{résilience} et une capacité à \highlight{influencer les autres} pour créer un environnement de travail positif et productif.

\section*{Valeur Ajoutée : Spécialiste Agile Multi-Squads}
\begin{itemize}[leftmargin=10pt]
    \item \highlight{Facilitation experte :} Supervise et facilite les cérémonies \highlight{Scrum} et \highlight{PI Planning} (priorisation, découpage, présentation des engagements)
    \item \highlight{Collaboration} optimale : Facilite et optimise la synergie entre squads et gère les \highlight{interdépendances}
    \item \highlight{Résolution de problèmes :} Diagnostique situations problématiques et met en place solutions durables
    \item \highlight{Coaching} et mentorat : Accompagne collègues moins expérimentés dans l'utilisation de la \highlight{méthodologie agile}
    \item \highlight{Gestion des risques :} Soulève risques et enjeux de manière constante et proactive
    \item \highlight{Communication} adaptée : Gère parties prenantes et maintient bonnes relations avec collaborateurs internes et externes
\end{itemize}


\section*{Compétences Clés}

\begin{minipage}[t]{0.55\textwidth}
\skilltitle{Frameworks Agile \& Facilitation}
\begin{itemize}[leftmargin=10pt]
    \item \highlight{Scrum Master certifié} : \highlight{Facilitation} de cérémonies et coaching \highlight{Agile}
    \item \highlight{SAFe (Scaled Agile Framework)} : Coordination multi-squads et \highlight{PI Planning}
    \item \highlight{Priorisation} et raffinement du backlog avec le Product Owner
    \item \highlight{Animation de réunion} et ateliers de travail
\end{itemize}

\skilltitle{Leadership \& Coaching}
\begin{itemize}[leftmargin=10pt]
    \item \highlight{Leadership} d'équipes (5+) avec un fort \highlight{sens de l'organisation}
    \item Coaching et mentorat des collègues dans l'agilité
    \item \highlight{Mobilisation} des équipes et capacité à \highlight{influencer les autres}
    \item \highlight{Gestion du stress} en environnement haute pression
\end{itemize}

\skilltitle{Delivery \& Qualité}
\begin{itemize}[leftmargin=10pt]
    \item \highlight{Amélioration continue} : Rétrospectives et optimisation des processus
    \item \highlight{Résolution de problèmes} : Diagnostic et levée d'obstacles
    \item \highlight{Gestion des risques} et enjeux de manière proactive
    \item Gestion des \highlight{interdépendances} entre squads pour accélérer la livraison
\end{itemize}
\end{minipage}
\hfill
\begin{minipage}[t]{0.42\textwidth}
\skilltitle{Outils \& Plateformes}
\begin{itemize}[leftmargin=10pt]
    \item \highlight{JIRA} (cert. Admin ACP-600), \highlight{Confluence}, suite Atlassian
    \item GitHub, Jenkins, Datadog, Splunk
    \item Automatisation tests (Selenium, Playwright, Robot...)
\end{itemize}

\skilltitle{Performance \& Pilotage}
\begin{itemize}[leftmargin=10pt]
    \item \highlight{Indicateurs clés de performance} (KPI) et tableaux de bord
    \item Suivi de la vélocité et rapports d'avancement
    \item \highlight{Communication} avec parties prenantes et reporting
    \item Gestion de budgets, échéanciers et actifs
\end{itemize}

\skilltitle{Soft Skills}
\begin{itemize}[leftmargin=10pt]
    \item \highlight{Communication} bilingue (français/anglais)
    \item \highlight{Écoute active} et empathie
    \item \highlight{Travail en équipe / Travail d'équipe} et \highlight{collaboration}
    \item \highlight{Initiative}, \highlight{résilience} et \highlight{orienté résultats}
\end{itemize}
\end{minipage}


\section*{Expérience Professionnelle}

\subsection*{Practice Lead Qualité | Responsable Pratique Qualité \& Facilitateur Agile} \hfill \textit{\color{mediumgray}Février 2025 -- Présent}\\
\textcolor{primaryblue}{\textbf{Banque Nationale du Canada (BNC)}} -- Montréal, QC
\begin{itemize}[leftmargin=22pt]
    \item \highlight{Facilitation Agile :} Supervise et facilite les cérémonies \highlight{Scrum} multi-squads démontrant d'excellentes compétences en \highlight{communication} et \highlight{écoute active}
    \item \highlight{Leadership} et \highlight{mobilisation :} Pilote la pratique qualité pan-train incluant équipes internationales en environnement \highlight{SAFe}
    \item \highlight{Coaching :} Encadre et développe une équipe de \keypoint{5+ professionnels QA} avec \highlight{initiative} et \highlight{sens de l'organisation}
    \item \highlight{PI Planning :} Prépare et participe activement aux comités de changement et \highlight{PI Planning} (priorisation, découpage, engagements)
    \item \highlight{Collaboration :} Optimise la synergie entre équipes Dev/QA/Ops et gère les \highlight{interdépendances}
    \item \highlight{Résolution de problèmes :} Implémente des \highlight{indicateurs clés de performance} et tableaux de bord (\highlight{JIRA}, Datadog, Splunk)
    \item \highlight{Gestion du stress :} Coordonne des POC en parallèle des opérations courantes, démontrant une forte \highlight{résilience}
    \item \highlight{Amélioration continue :} Organise un hackathon de test et promeut une culture d'apprentissage
    \item \highlight{Influencer les autres :} Agit comme ambassadeur des comportements BN (Agilité, Complicité, Pouvoir d'Agir)
\end{itemize}

\subsection*{Intégrateur Qualité Senior | Ingénieur Qualité Senior \& Chef de projet technique} \hfill \textit{\color{mediumgray}Janvier 2022 -- Février 2025}\\
\textcolor{primaryblue}{\textbf{Banque Nationale du Canada (BNC)}} via Groupe Astek -- Montréal, QC
\begin{itemize}[leftmargin=22pt]
    \item \highlight{Travail en équipe / Travail d'équipe :} Travaille étroitement avec les Product Owners pour définir des User Stories testables
    \item \highlight{Priorisation :} Soutient le PO dans la priorisation et le raffinement du backlog pour maximiser la valeur
    \item \highlight{Sens de l'organisation :} Conçoit la stratégie de test alignée avec la feuille de route Gestion de Patrimoine \& ECM
    \item \highlight{Communication :} Présente les métriques qualité aux parties prenantes pour assurer la transparence
    \item \highlight{Gestion des risques :} Identifie les risques techniques dès la conception, évitant des retards de livraison
    \item \highlight{Influencer les autres :} Conseille le PO sur la faisabilité technique et l'effort estimé
    \item \highlight{Initiative :} Crée des tableaux de bord qualité avec \highlight{indicateurs clés de performance}
\end{itemize}

\subsection*{QA Lead | Responsable Automatisation QA} \hfill \textit{\color{mediumgray}Juin 2021 -- Décembre 2021}\\
\textcolor{primaryblue}{\textbf{AODocs}} -- Paris, France
\begin{itemize}[leftmargin=22pt]
    \item \highlight{Orienté résultats :} Automatise \keypoint{100\% des tests de régression}, libérant \keypoint{20h/sprint} pour de nouvelles fonctionnalités
    \item \highlight{Initiative :} Évalue de nouveaux outils (tests mobiles Flutter, API Karate) et recommande leur adoption
    \item \highlight{Amélioration continue :} Implémente des design patterns (Page Object Model) adoptés par toute l'équipe
    \item \highlight{Communication :} Génère des rapports Allure facilitant la transparence sur l'état du produit
\end{itemize}


\subsection*{Ingénieur DevOps \& Coordinateur de projet technique} \hfill \textit{\color{mediumgray}Janvier 2020 -- Juin 2021}\\
\textcolor{primaryblue}{\textbf{ENGIE Digital}} via Groupe Hénix -- Paris, France
\begin{itemize}[leftmargin=22pt]
    \item \highlight{Sens de l'organisation :} Administre \highlight{JIRA} (\keypoint{8 projets, 50+ utilisateurs}), pilote la feuille de route et \highlight{priorise} les initiatives
    \item \highlight{Indicateurs clés de performance :} Met en place des tableaux de bord pour suivre la vélocité, les bugs et les performances
    \item \highlight{Orienté résultats :} Conçoit une chaîne CI/CD complète (Jenkins, Docker, Kubernetes) réduisant le \keypoint{time-to-market de 60\%}
    \item \highlight{Initiative :} Réalise des POC Robot Framework et Xray, recommande l'adoption basée sur le ROI
    \item \highlight{Collaboration :} Coordonne les équipes Dev, Ops et QA pour une intégration continue et gestion des \highlight{interdépendances}
\end{itemize}

\subsection*{Product Owner \& Responsable Delivery} \hfill \textit{\color{mediumgray}Juin 2019 -- Septembre 2019}\\
\textcolor{primaryblue}{\textbf{Hénix}} -- Montrouge, France
\begin{itemize}[leftmargin=22pt]
    \item \highlight{Priorisation :} Gère le backlog produit et priorise les fonctionnalités du plugin Xsquash (\highlight{JIRA} avec Squash TM)
    \item \highlight{Méthodologie agile :} Rédige et valide les histoires utilisateur pour l'intégration bi-directionnelle Jira/Squash
    \item \highlight{Écoute active :} Assure le support SaaS et collecte les retours utilisateurs pour une \highlight{amélioration continue}
    \item \highlight{Mobilisation :} Participe à la reprise du Centre de Service Thales (\keypoint{8500 utilisateurs}, 200+ projets)
\end{itemize}

\subsection*{Développeur logiciel | Développeur Full-Stack} \hfill \textit{\color{mediumgray}Mai 2017 -- Mai 2019}\\
\textcolor{primaryblue}{\textbf{MEDIAPOST}} via Groupe Hénix -- Paris, France
\begin{itemize}[leftmargin=22pt]
    \item Développe des webservices SOAP/REST (SOA BPEL) pour des applications de contractualisation et de facturation
    \item Crée une application web de gestion comptable (Node.js) avec import/export multi-formats (CSV, XML)
    \item \highlight{Initiative :} Automatise des processus métier avec Google Apps Script et macros Excel (gain \keypoint{15h/semaine})
    \item \highlight{Travail en équipe / Travail d'équipe :} Collabore avec l'équipe produit pour définir des spécifications techniques
\end{itemize}

\subsection*{Administrateur Plateformes \& Outils | Coordinateur technique} \hfill \textit{\color{mediumgray}Avril 2016 -- Avril 2017}\\
\textcolor{primaryblue}{\textbf{ENGIE IT}} via Groupe Hénix -- Saint-Ouen, France
\begin{itemize}[leftmargin=22pt]
    \item \highlight{Coordination :} Gère les environnements de production et coordonne les déploiements avec des équipes multiples
    \item Administre HP ALM Performance Center (\keypoint{150+ utilisateurs}) et assure le support technique
    \item \highlight{Amélioration continue :} Pilote les évolutions de la plateforme basées sur les besoins utilisateurs et le retour d'expérience
    \item \highlight{Résilience :} Maintient la disponibilité des services en environnement haute pression
\end{itemize}

\vspace{6pt}

\section*{Certifications \& Formations}

\textbf{PMP (en cours)} \hfill \textit{\color{mediumgray}2025}\\
\textbf{Scrum Fundamentals Certified (SFC)} \hfill \textit{\color{mediumgray}2024}\\
\textbf{Jira ACP-600 (Administration Jira)} \hfill \textit{\color{mediumgray}2019}\\
Autres : NeoLoad Pro, Selenium Professional, SoapUI, RPA \hfill \textit{\color{mediumgray}2019}

\section*{Formation Académique}

\textbf{Cursus Automaticien et DevOps} \hfill \textit{\color{mediumgray}2019}\\
AFCEPF, Montrouge, France -- Formation DevOps, CI/CD, Infrastructure as Code. 

\textbf{Master II (niveau) -- Architecture logicielle \& Analyste informaticien} \hfill \textit{\color{mediumgray}2017}\\
AFCEPF, France -- UML, Bases de données, Architectures distribuées, Big Data

\textbf{Licence en Gestion | Bachelor's Degree in Management} \hfill \textit{\color{mediumgray}2011}\\
École Supérieure de Gestion (E.S.G) -- Paris, France

\vspace{6pt}

\section*{Compétences Linguistiques}

\textbf{Français :} Langue maternelle \quad | \quad \textbf{Anglais :} Niveau avancé (Compétences professionnelles)\\
\textit{Capacité à collaborer étroitement avec des collègues à l'extérieur du Québec}

\end{document}
